% ISSUE 7 — Remark on Difficulty-Tier Counts: 24/27/20 -> 24/29/21 vs. "+3 Hard"
% Category: Text (pure text/arithmetic fix)
% Source: jmlr_paper_main.tex, subsec:difficulty-classification, l.635-662

\subsection{Difficulty Classification}\label{subsec:difficulty-classification}
\label{sec:bench_difficulty}

Tasks are categorised by difficulty into three tiers based on functional structure:

\begin{itemize}
  \item \textbf{Easy} ($n=24$): linear or simple rational relationships, such as
    Loan-to-Value, spot price from AMM reserve, collateral ratio, and simple staking APY.
    These tasks admit algebraic solutions with AST depth $\le 3$.

  \item \textbf{Medium} ($n=29$): nonlinear algebraic or exponential relationships,
    such as impermanent loss, Uniswap V3 liquidity range, constant-product price impact,
    and compounding staking returns.  These tasks require multi-step composition
    of algebraic operations.

  \item \textbf{Hard} ($n=21$): transcendental, probabilistic, or multivariate
    formulations.  This category includes Black--Scholes pricing (which requires
    the Gaussian CDF $\Phi$), all four Greeks, Portfolio Expected Shortfall for
    correlated assets, and multi-day VaR scaling.  These tasks are the primary
    source of LLM stochastic instability.
\end{itemize}

\begin{remark}
An earlier version of this benchmark comprised 71 tasks (Easy=24, Medium=27,
Hard=20).  The current benchmark (v3.0, 74 tasks) adds three additional hard
cases in multi-asset risk: correlated portfolio VaR, component ES, and
multi-collateral LTV.  All results in this paper refer to the 74-case v3.0 benchmark.
\end{remark}
